
\documentclass[english]{article}
%\documentclass[aps,pra,twocolumn,english]{revtex4}
%\documentclass[english]{iopart}
\usepackage[T1]{fontenc}
\usepackage[latin9]{inputenc}
\usepackage{babel}
\usepackage{graphicx,amsmath,amssymb,color}
\usepackage[normalem]{ulem}
\usepackage{amsfonts}
\usepackage[toc,page]{appendix}
%\usepackage[ocgcolorlinks,colorlinks=true,linkcolor=blue,citecolor=red]{hyperref}
\usepackage{hyperref}
\usepackage{latexsym}
\usepackage{amsfonts}
\usepackage{algpseudocode}
\usepackage{amsthm}
\usepackage{mathrsfs}
\usepackage{natbib}
\usepackage{color,verbatim}
\usepackage{psfrag}
\bibliographystyle{unsrt}
%\bibliographystyle{acm}
%\bibliographystyle{unsrtnat}
%\usepackage[style=numeric]{biblatex}

\usepackage{subcaption} % allows for subfigures


\usepackage[svgnames]{xcolor}
\usepackage{tikz}

\usepackage{algorithm}


%\input{tikzit/tikz_preamble.tex}  % load separately defined tikz preamble


\newcommand{\bra}[1]{\langle#1 |}
\newcommand{\ket}[1]{|#1 \rangle}
\newcommand{\braket}[2]{\left \langle #1 \mid #2 \right \rangle}
\newcommand{\bigbraket}[2]{\left \langle #1 \middle\vert #2 \right \rangle}
\newcommand{\ketbra}[2]{\vert #1 \rangle \! \langle #2 \vert}
\newcommand{\overlap}[2]{\left \langle #1 | #2 \right\rangle}
\newcommand{\average}[1]{\langle #1 \rangle}
\newcommand{\sandwich}[3]{\left \langle #1 \mid #2 \mid #3 \right\rangle}
\newcommand{\innerprod}[2]{\left \langle #1 | #2 \right\rangle}

\begin{document}


\title{How a bunch of quantum physics is actually the zeno effect in disguise}
\author{Nicholas Chancellor}


\maketitle



\today % added for drafting, remove in final version


{\small Disclaimer: lecture notes may contain ``typo'' type errors (factors of $2$, minus signs, etc...) if something like this looks wrong there is a good chance it is, please let me know}

\section{The basic Zeno effect}

\begin{itemize}
\item Classical (both in historical and physics sense) origin and name 
\begin{itemize}
\item Zeno of Elia 409-430 bc, ``arrow paradox''
\item Object seems still if observed in any one moment,  time is a compilation of moments, therefore motion is impossible if we think about ``observe'' all the individual moments
\item Arrow is fired but can't move because of this (spoiler/safety alert: doesn't actually work arrows fired at you can definitely kill you)
\item Does have a point that continuum is weird (actually weirder than anything quantum), but not the topic of this lecture
\end{itemize}
\item (Simplest) quantum version:
\begin{itemize}
\item Two level system, start in state $\ket{0}$ and undergo time evolution which would bring us to state $\ket{1}$ in time $\tau$ ($X$ is Pauli $X$ operation)
\begin{equation}
\ket{\psi(t)}=\exp\left(\frac{i t \pi }{2 \tau} X\right)\ket{0}=\cos(\frac{ t \pi }{2 \tau})\ket{0}+i \sin(\frac{ t \pi }{2 \tau})\ket{1}
\end{equation}
\item If we do nothing we end up in state $\ket{\psi(\tau)}=\ket{1}$
\item If we perform projective measurements we change the probabilities, in general, if we perform the evolution but measure $n$ times, (dividing equally into intervals of length $\frac{\tau}{n}$) the probability of still measuring $\ket{0}$ at every time is
\begin{equation}
P_{\ket{0}}=\cos^{2n}(\frac{  \pi }{2 n})=(1-\sin^2(\frac{ \pi }{2n}))^n
\end{equation}
\item If we assume $n$ is large, we can Taylor expand, and then use a well known infinite series expansion for the exponential $\lim_{n\rightarrow \infty}(1-\frac{x}{n})^n=\exp(-x)$
\begin{equation}
P_{\ket{0}}\approx(1-(\frac{ \pi }{n})^2)^n\approx \exp(-\frac{ \pi ^2}{n})
\end{equation}
\item Clearly as $n\rightarrow \infty$, $P_{\ket{0}} \rightarrow 1$
\end{itemize}
\item What does the simplest quantum version need?
\begin{itemize}
\item Projective measurement, otherwise $P_{\ket{0}}$ would be unaffected by measurement
\item Dependence on amplitude squared, in a fictional version of QM where probability scales as amplitude rather than amplitude squared, we would get $P_{\ket{0}}\approx \exp(- \pi)$, tend toward fixed value less than $1$
\end{itemize}
\end{itemize}

\section{Extensions, similar effects}


\begin{itemize}
\item Don't measure, but make levels distinguishable with a phase, in other words, apply Pauli $Z$ at $n$ evenly spaced instances
\begin{itemize}
\item Even $n=1$, the dynamics will reverse perfectly since $X$ and $Z$ anti-commute ($XZ=-ZX$), therefore  
\begin{equation}
\exp\left(\frac{i \pi }{4} X\right)Z\exp\left(\frac{i \pi }{4} X\right)\ket{0}=Z\exp\left(\frac{-i \pi }{4} X\right)\exp\left(\frac{i \pi }{4} X\right)\ket{0}=\ket{0}
\end{equation}
\item Won't work so neatly in more complicated systems, but cancellation will still happen at large $n$
\end{itemize}
\item Continuously apply a phase between energy levels to stop evolution $X \rightarrow X-\frac{\Delta}{2} Z$ (note $\Delta$ plays an equivalent role to $n$ but doesn't need to be an integer)
\begin{itemize}
\item How do we represent $\exp\left(\frac{i t \pi }{2 \tau} (X-\frac{\Delta}{2} Z)\right)$?
\item $X+Z$ is a 2x2 matrix, we can diagonalize it by hand (exponentiation is trivial in diagonal basis)
\begin{equation}
\det\left(\begin{array}{cc} -\frac{\Delta}{2}-\lambda & 1 \\ 1 & \frac{\Delta}{2}-\lambda \end{array}\right)=0
\end{equation}
\item $(\frac{\Delta}{2}-\lambda)(-\frac{\Delta}{2}-\lambda)-1=0$ expanding and rearranging $\lambda^2=1+\frac{\Delta^2}{4}$
\item $\lambda_{\pm}=\pm\sqrt{1+\frac{\Delta^2}{4}}$
\item Eigenvector equation $-\frac{\Delta}{2}v_0+v_1=\lambda_{\pm}v_0$
\begin{equation}
\frac{v_1}{v_0}=\frac{\Delta}{2}\pm\sqrt{1+\frac{\Delta^2}{4}}
\end{equation}
\item Taking $\lambda_{-}$, we see that $\frac{v_1}{v_0}\rightarrow 0$ as $\Delta \rightarrow \infty$
\item This implies that for large $\Delta$, $\ket{0}$ is an eigenstate of $X-\frac{\Delta}{2} Z$, since exponential of matrix is diagonal in same basis, it is also an eigenstate of $\exp\left(\frac{i t \pi }{2 \tau} (X-\frac{\Delta}{2} Z)\right)$ for any $t$, therefore never leaves $\ket{0}$
\end{itemize}
\item The previous example was just the adiabatic theorem in disguise
\begin{itemize}
\item Consider energy eigenbasis of Hamiltonian, and only first two energy levels, time dependent energy difference $\frac{\Delta}{2}Z$
\item Basis is rotating over time, transition strength will (up to an irrelevant phase) be equal to $\left|\braket{\dot{\phi}_0}{\phi_1}\right|$, where $\phi_0$ and $\phi_1$ are the ground and first excited state and the dot indicates time derivative
\item Since the $X$ and $Z$ terms are time dependent and non-commuting, need path integral to represent evolution
\begin{equation}
\ket{\psi(\tau)}=\int_0^\tau Dt \exp\left(\frac{\Delta(\frac{t}{\tau})}{2}Z+\left|\braket{\dot{\phi}_0(\frac{t}{\tau})}{\phi_1(\frac{t}{\tau})}\right|X\right)\ket{0}
\end{equation}
\item By previous argument as long as $\frac{\Delta(\frac{t}{\tau})}{2}\gg\left|\braket{\dot{\phi}_0(\frac{t}{\tau})}{\phi_1(\frac{t}{\tau})}\right|$, than time evolution than $\ket{0}$ will be an eigenstate of the time evolution at all times and the system will remain in it's ground state
\item $\left|\braket{\dot{\phi}_0(\frac{t}{\tau})}{\phi_1(\frac{t}{\tau})}\right| \propto \frac{1}{\tau}$ for fixed $\frac{t}{\tau}$, therefore as long as $\Delta(\frac{t}{\tau})>0$ it will always be possible to pick a large enough $\tau$ so that the ground state is maintained with arbitrarily high probability
\end{itemize}
\item What if instead of phases we have particle loss from the $\ket{1}$ state?  
\begin{itemize}
\item Simplest version, add complex term to time evolution evolution with $\exp\left(\frac{i t \pi }{2 \tau} (X+i\Gamma\ketbra{1}{1})\right)$
\item Same trick, diagonalize matrix by hand 
\begin{equation}
\det\left(\begin{array}{cc} -\lambda & 1 \\ 1 & i\Gamma-\lambda \end{array}\right)=0
\end{equation}
\item $\lambda^2-i\Gamma\lambda-1=0$ from an application of the quadratic formula $\lambda_{\pm}=\frac{i\Gamma \pm\sqrt{-\Gamma^2+4}}{2}$
\item Simple eigenvector equation $v_1=\lambda_{\pm}v_0$
\item  Since $\lambda_{-}\rightarrow 0$, than $\frac{v_1}{v_0}\rightarrow 0$ for this eigenvalue and therefore $\ket{0}$ tends toward an eigenstate, just as it did with phases
\end{itemize}
\item Beyond $2$ level Hermitian matrix
\begin{itemize}
\item Except for special cases we cannot easily diagonalize matrices by hand, but...
\item Let's say I have a matrix with a submatrix $M'$ with very large eigenvalues such that $\lambda_{\mathrm{min}}\le|\lambda_{M',i}| \quad \forall i$
\item Let us say that the sum of the magnitudes all matrix elements outside of $M'$ obeys
\begin{equation}
\sum_{i,j\notin M'}|A_{i,j}|\ll\lambda_{\mathrm{min}}
\end{equation}
\item Since this sum of the magnitudes cannot contribute more than $\sum_{i,j\notin M'}|A_{i,j}|$ to the eigenvalues, it follows that in the limit $\frac{\sum_{i,j\notin M'}|A_{i,j}|}{\lambda_{\mathrm{min}}}\rightarrow 0$ $\lambda_{M',i}$  will be unchanged
\item Since adding support to the eigenvectors outside of $M'$ would reduce the eigenvalues, it implies that the associated eigenvectors will be unchanged
\item Since eigenvectors of a Hermitian matrix are orthogonal, this implies that all other eigenvectors will only have support outside of $M'$
\item Therefore, many kinds of strong interactions can lead to a Zeno-effect-like separation of parts of the system with fast dynamics and those with slow dynamics
\end{itemize}
\section{Non-trivial subsystems}
\item Zeno effect doesn't stop \textbf{all} dynamics, it only prevents dynamics from entering certain states
\begin{itemize}
\item Consider $2$ qubits, use Zeno effect (by strong decay) to stop system from entering $\ket{11}$ state
\item Start in separable state $\ket{+0}=\frac{1}{\sqrt{2}}(\ket{00}+\ket{10})$
\item Apply $X$ to second qubit (n.b.~cannot enter $\ket{11}$): $\frac{1}{\sqrt{2}}(\ket{01}+\ket{10})$ \textbf{entangled} state
\item Zeno effect can turn separable states + separable operations into entangled states
\end{itemize}
\item We can make quantum gates with this... but need three qubits because both need to start in arbitrary states
\begin{itemize}
\item Simple example: use Zeno effect to prevent system from entering $\ket{111}$ state 
\item Start in $\ket{\psi_{01}}\ket{0}=(a_{00}\ket{00}+a_{01}\ket{01}+a_{10}\ket{10}+a_{11}\ket{11})\ket{0}$, last qubit guarantees system not in $\ket{111}$
\item Apply $X$ to last qubit $\ket{\psi}\ket{0} \rightarrow (a_{00}\ket{00}+a_{01}\ket{01}+a_{10}\ket{10})\ket{1}+a_{11}\ket{11}\ket{0}$
\item Apply $Z$ to last qubit again yielding $-(a_{00}\ket{00}+a_{01}\ket{01}+a_{10}\ket{10})\ket{1}+a_{11}\ket{11}\ket{0}$
\item Apply another $X$ to last qubit $(a_{00}\ket{00}-a_{01}\ket{01}-a_{10}\ket{10}+a_{11}\ket{11})\ket{0}$
\item Total operation $\ket{\psi}\ket{0} \rightarrow -(\mathrm{CZ}\ket{\psi})\ket{0}$, we have applied a controlled $Z$ on first two qubits up to a global phase
\end{itemize}
\item What does an adiabatic analog of this look like? Geometric gates
\begin{itemize}
\item Same idea, cannot leave the ground state manifold
\item Geometric picture, moving in closed loops can cause rotations within ground state manifold (Holonomic quantum computing)
\item Generally complicated, a relatively simple example:
\item Three level system $\ket{0},\ket{1},\ket{e}$, initially $\ket{0}$ and $\ket{1}$ are degenerate, $\ket{e}$ has a higher energy to start with $H_0=\ketbra{e}{e}$
\begin{enumerate}
\item Adiabatically evolve for time $\tau$ with
\begin{equation}
H_1(t)=\left(\sqrt{1-\frac{t}{\tau}}\ket{e}+\sqrt{\frac{t}{\tau}}\ket{1}\right)\left(\sqrt{1-\frac{t}{\tau}}\bra{e}+\sqrt{\frac{t}{\tau}}\bra{1}\right)
\end{equation}
 ground state manifold is now $\{\ket{0},\ket{e}\}$, where $\ket{0}\rightarrow \ket{0}$ and $\ket{1}\rightarrow \ket{e}$
 \item Adiabatically evolve for time $\tau$ with
\begin{equation}
H_2(t)=\left(\sqrt{1-\frac{t}{\tau}}\ket{1}+\sqrt{\frac{t}{\tau}}\ket{0}\right)\left(\sqrt{1-\frac{t}{\tau}}\bra{1}+\sqrt{\frac{t}{\tau}}\bra{0}\right)
\end{equation}
 ground state manifold is now $\{\ket{1},\ket{e}\}$, where $\ket{0}\rightarrow \ket{0}\rightarrow \ket{1}$ and $\ket{1}\rightarrow \ket{e} \rightarrow \ket{e}$
 \item Adiabatically evolve for time $\tau$ with
\begin{equation}
H_3(t)=\left(\sqrt{1-\frac{t}{\tau}}\ket{0}+\sqrt{\frac{t}{\tau}}\ket{e}\right)\left(\sqrt{1-\frac{t}{\tau}}\bra{0}+\sqrt{\frac{t}{\tau}}\bra{e}\right)
\end{equation}
 ground state manifold is now $\{\ket{1},\ket{0}\}$, where $\ket{0}\rightarrow \ket{0}\rightarrow \ket{1}\rightarrow \ket{1}$ and $\ket{1}\rightarrow \ket{e} \rightarrow \ket{e}\rightarrow \ket{0}$, moreover $H_3(\tau)=H_0$
\end{enumerate}
\item We have flipped a spin purely by geometry, we adiabatically traversed a loop in phase space and the enclosed curvature flipped a bit!
\item Note that as long as $\tau$ is big enough, the action of this gate does not depend on the value of $\tau$
\end{itemize}
\end{itemize}





\end{document}